\section{The cases of four- and five-cycles}
\label{sec:small}

We now rule out the minimal counterexample when $k=3$ or $k=4$.
Only the triangle guaranteed by Lemma~\ref{clq:triangle} is needed
from the preceding section. The argument examines the neighbors
outside a clique and uses the absence of the required exact cycle
to separate their neighborhoods.

\begin{proposition}\label{small:cases}
Let $k\in\{3,4\}$ and let $a$ be an integer with $2\le a\le k$.
There is no graph $G$ with
\[
  |V(G)|=ka+1,\qquad \alpha(G)\le a,
  \qquad G\text{ contains no }C_{k+1},
\]
such that $|N_G[I]|\ge k|I|+1$ for every nonempty independent set
$I\subseteq V(G)$.
\end{proposition}

\begin{proof}
Suppose such a graph exists. Applying the neighborhood bound to a
singleton gives $\delta(G)\ge k$. Lemma~\ref{clq:triangle} supplies
a triangle, and a clique of order $k+1$ would contain the forbidden
cycle, so $3\le\omega(G)\le k$.

For a clique $Q=\{q_1,\ldots,q_t\}$, let $F=G-Q$ and put
\[
  U_i=N_G(q_i)\cap V(F)\qquad(1\le i\le t).
\]
We will use this notation with different choices of $Q$, specifying
the choice each time. A path between two distinct vertices of $Q$
whose internal vertices lie in $F$ can be completed using edges of
$Q$. The exact numbers of internal vertices that are forbidden are
checked below.

\medskip
\noindent\textbf{The case $k=3$.}
Choose a triangle $Q=\{q_1,q_2,q_3\}$. Each $U_i$ is nonempty,
since $q_i$ has two neighbors inside $Q$ and degree at least three.
For distinct $i,j$, write $h$ for the remaining index. A common
vertex $u\in U_i\cap U_j$ would give the cycle
\[
  q_i,u,q_j,q_h,q_i
\]
of order four. If $u\in U_i$ and $v\in U_j$ were adjacent, then
\[
  q_i,u,v,q_j,q_i
\]
would be a cycle of order four. Thus the sets $U_1,U_2,U_3$ are
pairwise disjoint and pairwise anticomplete.

Choosing one vertex from each $U_i$ gives an independent set of
order three. Hence $a=3$. If any $U_i$ contained two nonadjacent
vertices, adjoining one vertex from each of the other two sets would
give an independent set of order four. Therefore each $U_i$ is a
clique. Moreover, $\{q_i\}\cup U_i$ is a clique and there is no
$K_4$, so
\[
  1\le |U_i|\le2.
\]
Since $|V(F)|=7$, the set
\[
  W=V(F)\setminus(U_1\cup U_2\cup U_3)
\]
is nonempty. Every vertex of $W$ has no neighbor in $Q$ by the
definition of the $U_i$.

Fix $w\in W$. If $w$ saw two distinct vertices $u,v\in U_i$, then
$w,u,q_i,v,w$ would be a $C_4$. Thus $w$ has at most one neighbor
in each $U_i$. On the other hand, if $w$ missed a vertex $u\in U_i$
and a vertex $v\in U_j$ for distinct $i,j$, then
$\{w,u,v,q_h\}$ would be an independent set of order four. Indeed,
$w$ misses $Q$, the two exterior sets are anticomplete, and $q_h$
has no neighbor in either of them. It follows that $w$ is adjacent
to every vertex of at least two of the $U_i$. These two sets must
therefore be singletons.

There are consequently at least two singleton sets among the three
$U_i$. If exactly two are singletons, the three sizes are $1,1,2$,
so $|W|=3$. Every vertex of $W$ sees the two singleton vertices,
say $u,v$. Distinct $w,w'\in W$ then give the cycle
$w,u,w',v,w$ of order four. If all three $U_i$ are singletons,
then $|W|=4$. Each $w\in W$ sees at least two of the three singleton
vertices. Choose one pair of such neighbors for each $w$. There are
only three pairs, so two vertices $w,w'$ choose the same pair
$\{u,v\}$, again giving $w,u,w',v,w$. This rules out $k=3$.

\medskip
\noindent\textbf{The case $k=4$ with a clique of order four.}
Suppose first that $Q=\{q_1,q_2,q_3,q_4\}$ is a clique. A path
between distinct $q_i,q_j$ with one, two, or three internal vertices
outside $Q$ can be completed to a $C_5$ by a path in $Q$ of length
three, two, or one, respectively. More explicitly, with $q_h,q_\ell$
the other clique vertices, the prohibited configurations yield
\[
\begin{gathered}
  q_i,u,q_j,q_h,q_\ell,q_i,\\
  q_i,u,v,q_j,q_h,q_i,\\
  q_i,u,z,v,q_j,q_i.
\end{gathered}
\]
In every row the exterior vertices are distinct and lie outside
$Q$, so the displayed cycle has exactly five vertices.

The sets $U_i$ are nonempty. Their closed neighborhoods in $F$ are
pairwise disjoint: an intersection of $N_F[U_i]$ and $N_F[U_j]$
would give vertices $u\in U_i$, $v\in U_j$ at distance at most two
in $F$, with $u=v$ allowed. A shortest such path, joined to $q_i$
and $q_j$, is one of the three prohibited configurations. In
particular, every vertex $u\in U_i$ has exactly one neighbor in
$Q$, so it has at least three neighbors in $F$. It follows that
\[
  |V(F)|\ge\sum_{i=1}^4|N_F[U_i]|\ge4\cdot4=16.
\]
But $|V(F)|=4a-3\le13$, a contradiction. We may therefore assume
that $\omega(G)=3$.

\medskip
\noindent\textbf{Triangles with disjoint exterior neighbor sets.}
For any triangle $Q=\{q_1,q_2,q_3\}$, every $U_i$ has size at
least two. A path between distinct triangle vertices with two or
three internal vertices outside $Q$ gives a $C_5$. In particular,
if distinct $u\in U_i$, $v\in U_j$ are adjacent, or are joined by
a two-edge path $u,z,v$ in $F$, the respective cycles are
\begin{equation}
  q_i,u,v,q_j,q_h,q_i
  \qquad\text{and}\qquad
  q_i,u,z,v,q_j,q_i,
  \label{small:five-cycles}
\end{equation}
where $h$ is the remaining triangle index.

Suppose for the moment that $U_1,U_2,U_3$ are disjoint. By
\eqref{small:five-cycles}, their closed neighborhoods in $F$ are
also pairwise disjoint. A maximum independent subset $J$ of $U_i$
has exactly one neighbor in $Q$, namely $q_i$. Thus the assumed
expansion bound gives
\begin{equation}
  |N_F[U_i]|\ge |N_F[J]|
  =|N_G[J]|-1\ge4\alpha(U_i).
  \label{small:one-neighbor-count}
\end{equation}
Here $|V(F)|=4a-2\le14$. If any $U_i$ had independence number at
least two, the three disjoint neighborhoods would together have
at least $8+4+4=16$ vertices. Hence every $U_i$ is a clique. Its
size is exactly two, because it has size at least two and adjoining
$q_i$ cannot produce a $K_4$.

If the exterior neighbor sets were disjoint for every triangle,
write $U_i=\{u_i,v_i\}$ and apply that assertion to the triangle
$T_i=\{q_i,u_i,v_i\}$. Each of $u_i,v_i$ has at least two neighbors
outside $T_i$, and the two sets of such neighbors are disjoint.
All of these vertices lie in $F-U_i$: the original disjointness
prevents $u_i$ or $v_i$ from seeing a vertex of $Q\setminus\{q_i\}$.
Consequently $|N_F[U_i]|\ge2+2+2=6$ for every $i$. The three
original closed neighborhoods are disjoint, so this would require
at least $18$ vertices in $F$, again impossible. Some triangle
therefore has overlapping exterior neighbor sets.

\medskip
\noindent\textbf{The necessary structure at an overlapping triangle.}
Take a triangle for which
$I=U_1\cap U_2\ne\varnothing$. No other pair of exterior sets
overlaps. Indeed, a vertex in all three sets would form a $K_4$
with $Q$. If $x\in U_1\cap U_2$ and
$y\in U_1\cap U_3$ are distinct, then
\[
  q_2,x,q_1,y,q_3,q_2
\]
is a $C_5$; if instead $y\in U_2\cap U_3$, use
$q_1,x,q_2,y,q_3,q_1$. Also $I$ is independent, since adjacent
vertices $x,y\in I$ would give the $K_4$ on $q_1,q_2,x,y$.

The four sets
\[
  I,\qquad A=U_1\setminus I,\qquad
  B=U_2\setminus I,\qquad D=U_3
\]
have pairwise disjoint closed neighborhoods in $F$. To check this
even for pairs involving $I$, choose distinct triangle neighbors
for vertices in the six possible pairs of classes as follows:
\[
\begin{array}{c|cccccc}
\text{classes}&(I,A)&(I,B)&(I,D)&(A,B)&(A,D)&(B,D)\\ \hline
\text{neighbors}&(q_2,q_1)&(q_1,q_2)&(q_1,q_3)&
                  (q_1,q_2)&(q_1,q_3)&(q_2,q_3).
\end{array}
\]
If two vertices from the indicated classes were adjacent or shared
a neighbor in $F$, the corresponding distinct triangle neighbors
would close one of the cycles in \eqref{small:five-cycles}. Since
the classes themselves are disjoint, their closed neighborhoods
are disjoint as claimed.

Write $r=|I|$. The set $I$ is independent and its only neighbors in
$Q$ are $q_1,q_2$. Hence
\begin{equation}
  |N_F[I]|=|N_G[I]|-2\ge4r-1.
  \label{small:overlap-count}
\end{equation}
Each nonempty one of $A,B,D$ has exactly one triangle neighbor,
so the argument of \eqref{small:one-neighbor-count} bounds the
size of its closed neighborhood below by four times its
independence number. These bounds may be added because the four
closed neighborhoods are disjoint.

If $r=1$, then both $A$ and $B$ are nonempty, since
$|U_1|,|U_2|\ge2$, and $D$ is nonempty as well. The four
neighborhoods would have total size at least
$3+4+4+4=15>14$. If $r\ge3$, the neighborhoods of $I$ and $D$
alone would have total size at least $11+4=15>14$. Thus $r=2$.
In this case the neighborhoods of $I$ and $D$ already contribute
at least $7+4=11$ vertices; a nonempty $A$ or $B$ would increase
the bound to $15$. Therefore
\begin{equation}
  |I|=2,\qquad U_1=U_2=I.
  \label{small:overlap-structure}
\end{equation}
Moreover, $\alpha(D)\ge2$ would give $7+8=15$ vertices, so $D$
is a clique. Its size is exactly two, by the degree bound and the
absence of a $K_4$.

This argument proves a necessary statement for \emph{every}
triangle with overlapping exterior neighbor sets: exactly one
pair overlaps, both sets in that pair equal an independent set
of size two, and the third exterior neighbor set is a clique of
size two. No condition on the exterior sets of other triangles
was used in proving this statement.

\medskip
\noindent\textbf{The final contradiction.}
Write $I=\{w,z\}$ and consider the triangle
$T=\{q_1,q_2,w\}$. By \eqref{small:overlap-structure}, the
exterior neighbor sets of its first two vertices are both
\[
  N_G(q_1)\setminus T=N_G(q_2)\setminus T=\{q_3,z\}.
\]
The necessary statement just proved therefore applies to $T$.
Its third vertex $w$ has exactly two neighbors outside $T$.
Neither $q_3$ nor $z$ is adjacent to $w$: adjacency to $q_3$
would give a $K_4$ with the original triangle, and $I$ is
independent. Thus the neighbors of $w$ outside $T$ are precisely
its neighbors in the original graph $F$, and $d_F(w)=2$.
The same reasoning applies to $z$. Consequently,
\[
  |N_F[I]|\le |I|+d_F(w)+d_F(z)=2+2+2=6,
\]
whereas \eqref{small:overlap-count} gives $|N_F[I]|\ge7$.
This contradiction completes the case $k=4$ and the proof.
\end{proof}

By Lemma~\ref{red:expansion}, the minimal counterexample satisfies
all the assumptions of Proposition~\ref{small:cases}. Hence we
may assume $k\ge5$ from now on.
